\section{Introduction}

The previous chapter described the PatchOps architecture as a structured repair pipeline. This chapter maps that
design onto the actual implementation. The goal is not to list every file mechanically, but to explain how the
main modules cooperate to realize repository-aware, safety-gated, abstention-aware CI repair in code.

The implementation centers on the \texttt{patchops.engine} package, with supporting validation, adapter, and
benchmark code around it. This separation is intentional. The engine handles reasoning and decision logic. The
validation layer handles CI-style re-execution. The benchmark scripts use the engine as a reusable subject of
study rather than duplicating repair logic outside the system.

\section{Implementation Structure}

The PatchOps implementation is organized around a small set of coherent modules rather than a monolithic agent
script. At the highest level, the codebase is divided into:

\begin{enumerate}
\item \textbf{engine modules} under \path{patchops/engine/}, which implement diagnosis, retrieval, planning,
generation, safety, and final decision logic,
\item \textbf{validation modules} under \path{patchops/validation/}, which provide local CI-style
re-execution,
\item \textbf{adapters} under \path{patchops/adapters/}, which support GitHub- and CLI-facing integration,
\item \textbf{benchmark scripts} under \path{benchmarks/}, which evaluate the engine under controlled
experimental protocols, and
\item \textbf{acceptance and verification tests} under \path{tests/}, which enforce regression guarantees on
the core behavior.
\end{enumerate}

This layered structure is valuable for a thesis artifact because it keeps the experimental system reusable. The
same \texttt{diagnose(...)} entry point used in the live or local workflow can also be exercised by the
benchmark harness, making the evaluation more faithful to the actual system behavior.

\section{The Engine Entry Point}

The practical center of the implementation is the decision routine in \path{patchops/engine/decide.py}. This
module realizes the state machine described earlier: log extraction, classification, flaky detection, context
retrieval, root-cause analysis, intended-action planning, semantic ambiguity downgrading, patch generation,
safety gating, plausibility scoring, and final decision selection.

The implementation returns a structured \texttt{Result} object rather than plain text. This object bundles the
evidence, classification, root cause, decision, optional patch, optional safety report, and the reasoning
trace. Supporting dataclasses in \texttt{types.py} make the pipeline explicit and serialization-friendly.
Architecturally, this means every stage writes to a shared, typed interface rather than leaking unstructured
prompt outputs across the system.

\section{Evidence Compression and Classification}

The first implementation layer is evidence compression. The log extractor in
\path{patchops/engine/log_extractor.py} performs deterministic preprocessing: it strips ANSI noise, removes
timestamp clutter, detects known signatures, finds the failed step, and isolates a critical error window. This
keeps the later prompts short and makes the engine usable even without a live model.

Classification is implemented in \texttt{classifier.py}. Its policy is hybrid. Strong deterministic signatures
remain useful, but noisy real logs often require model adjudication. When an LLM is available, the classifier
uses the rule-based hint and the learned-classifier hint as inputs to an LLM-primary decision. When no LLM is
available, a confident learned-classifier prediction becomes the offline primary; otherwise rule signatures are
used as the fallback. This preserves the thesis claim that PatchOps degrades gracefully rather than collapsing
without API access.

\section{Provider-Agnostic LLM Layer}

One of the most important implementation choices is the provider-agnostic LLM abstraction in
\path{patchops/engine/llm.py}. Instead of hard-wiring the system to one vendor, PatchOps supports OpenAI,
Anthropic, and OpenAI-compatible providers behind a shared interface. The implementation also supports local or
alternative OpenAI-style endpoints through a configurable base URL. If no provider is configured, the engine
returns cleanly to deterministic behavior.

This design serves three purposes. First, it makes experiments reproducible by pinning model roles through
configuration. Second, it lets the system distinguish between strong and cheaper models for different stages,
such as classification versus patch generation. Third, it turns LLM access into an optional acceleration layer
rather than a hard dependency. That is especially important for a thesis claiming reproducibility and careful
systems design rather than merely cloud-backed prompt orchestration.

\section{Repository Context and Localization}

Repository awareness is implemented through a pair of cooperating modules:
\path{context_retriever.py} and \path{localizer.py}, with \path{repo_index.py} extending the design for
full-repository search. The context retriever collects category-relevant files such as workflow YAML,
manifests, lockfiles, configuration files, and evidence-named source files. It does so under size caps to keep
prompt cost bounded.

The localizer makes file selection an explicit, measurable subproblem. It implements multiple strategies,
including log-path ranking, category-aware reranking, LLM candidate generation, hybrid fusion, deterministic
full-repository BM25 ranking, and semantic embedding reranking. The repository index in
\path{repo_index.py} builds a candidate universe from the repository clone itself, applies path-aware BM25,
and optionally reranks the deterministic pool. This modular structure is important because it lets the thesis
study localization as its own bottleneck rather than treating file selection as an invisible side effect of a
single large prompt.

\section{Root-Cause Analysis and Patch Planning}

Once context has been gathered, PatchOps performs root-cause analysis through \path{rca_agent.py}. This
module combines rule-based templates with optional LLM refinement. The implementation reflects the thesis
philosophy of bounded intelligence: rules provide a defensible baseline and zero-cost fallback, while the model
can produce a more contextual summary and affected-file hypothesis when available.

The planner in \path{patch_planner.py} then turns diagnosis into intended action. It maps category,
confidence, and flakiness into the action policy described earlier, using pinned thresholds and fixed risk
classes from \path{config.py}. The semantic ambiguity guard is also implemented here, allowing the engine to
downgrade an intended patch to preview when type-checking or logic-oriented signals suggest that human judgement
is still required.

The learned decision gate is implemented in \path{decision_gate.py} and wired into
\path{decide.py}. It is opt-in and replaces only the low-risk high-confidence act-versus-hold test when the
environment enables it and the model artifact is available. It does not bypass flaky rerun routing, high-risk
escalation, the safety gate, or the plausibility gate. This narrow integration is deliberate: the learned gate
removes the old magic threshold on the low-risk lane while preserving deterministic safety boundaries.

\section{Patch Generation}

Patch generation is implemented in \path{patch_generator.py}. The module contains both deterministic repair
paths and LLM-backed patch generation. The deterministic path handles important mechanical cases such as runtime
version mismatches, where the repair can be computed directly from repository and log evidence. The LLM path
handles a broader set of repairable categories by asking the model to return full edited file contents and then
computing the unified diff locally.

This is an important implementation detail. Asking the model for raw diff text directly is brittle. By asking
for full edited file contents and deriving the diff in code, PatchOps keeps patch formation more controllable
and more likely to remain \texttt{git apply}-compatible. The generator also produces review-oriented PR
metadata, including a title and a security-conscious PR body, reinforcing that the output is meant for
inspection rather than blind execution.

\section{Safety and Plausibility Controls}

The two most thesis-critical safety modules are \path{safety_gate.py} and
\path{patch_plausibility.py}. The safety gate implements deterministic checks over the candidate patch. It
blocks forbidden paths, category-level allow-list violations, secret introduction, excessive patch breadth, and
reward-hacking patterns. Because this logic is deterministic and explicit, it can be measured directly in the
evaluation rather than inferred informally.

The plausibility module implements a second layer of control. It estimates whether a generated patch is likely
to be on-target using only inference-time signals such as evidence locality, file-type fit, minimality, and
honesty. Low-plausibility patches are downgraded from automatic patch to preview. Together, these two modules
operationalize the thesis argument that CI-green and patch-generated are not strong enough success conditions on
their own.

\section{Learned Classifier Integration}

The learned classifier is implemented as an additive layer in \path{ml_classifier.py}. It loads a trained
artifact produced by the benchmark training script and predicts category plus confidence from the same compressed
evidence representation used at inference time. If the artifact is missing or the supporting dependencies are
unavailable, the module simply yields no prediction and the engine falls back to its earlier rule and LLM path.

This is an important implementation property. The classifier does not destabilize the rest of the system. It is
integrated as a graceful enhancement for reproducibility and low-cost inference, not as a mandatory dependency
that would change the rest of the engine's control flow.

The classifier artifact also carries the shipped isotonic calibrator used by later decision stages. When present,
\path{ml_classifier.py} applies this calibrator before emitting confidence; when absent, it falls back to raw
confidence. This makes calibrated confidence an inference-time engine property rather than only an offline
plotting result.

\section{Validation Harness}

The CI validation layer is implemented under \path{patchops/validation/}. Its local runner parses common
GitHub Actions workflow shapes, extracts runnable shell steps, and executes them either in Docker or on the
host as a fallback. The implementation is deliberately modest about scope: it supports common setup-node and
setup-python patterns plus shell-based workflow steps, but it is not a full GitHub Actions emulator.

This bounded implementation is nevertheless valuable. It gives PatchOps a local, repeatable mechanism for
testing whether a candidate patch changes the observable CI outcome. The validation layer is therefore what
connects the abstract notion of repair to the measurable Pass@1 and Pass@3 style metrics used later in the
thesis.

\section{Benchmarks as First-Class Implementation Clients}

The benchmark scripts are not bolted on as an afterthought. They are designed as disciplined clients of the same
engine. For example, \path{benchmarks/run_eval.py} drives the diagnosis and patch pipeline across ablation
variants, \path{run_localization.py} isolates file-selection strategies, and the later benchmark scripts
evaluate full-repository retrieval, the learned classifier, selective prediction, plausibility, reward hacking,
agentic localization, and the learned decision gate.

This matters because it means the measured system is the real system. The thesis is not evaluating simplified
mock-ups of PatchOps stages; it is evaluating the actual implemented modules under controlled harnesses.

\section{Acceptance and Regression Tests}

PatchOps also includes acceptance and regression tests that enforce the system's core contract. The acceptance
harness checks end-to-end diagnosis on the ten-case example set, verifies that the flagship deterministic patch
applies cleanly, confirms correct rerun and escalate behavior on flaky and secret-related examples, requires a
reasoning trace, and checks both the semantic ambiguity and reward-hacking protections.

These tests are important in a thesis context because they distinguish two levels of confidence. Experimental
results tell us how the system behaves on a benchmark. Acceptance checks tell us whether the intended design
properties still hold after implementation changes. Together, they strengthen the system as a research artifact.

\section{Implementation Principles}

Several implementation principles recur across the codebase:

\begin{enumerate}
\item \textbf{Graceful degradation.} Optional LLM, classifier, retrieval, and decision-gate enhancements are
additive rather than destructive; the system retains deterministic functionality when they are unavailable.
\item \textbf{Bounded cost.} File size caps, prompt truncation, top-k limits, and staged retrieval keep model
usage and runtime manageable.
\item \textbf{Deterministic policy barriers.} Safety, ambiguity, and risk checks are enforced in code rather
than left to prompt compliance.
\item \textbf{Modularity.} Retrieval, classification, generation, safety, and validation are separable enough to
be ablated experimentally.
\item \textbf{Traceability.} Each stage contributes explicit outputs to a reasoning trace that can be inspected
later by users or benchmark harnesses.
\end{enumerate}

These principles are what make PatchOps more than a prompt chain. They turn it into a controlled software
system whose behavior can be studied systematically.

\section{Chapter Conclusion}

This chapter mapped the PatchOps architecture to its concrete implementation. The engine centers on a typed,
traceable decision pipeline; the LLM layer is provider-agnostic and optional; repository retrieval and
localization are modular and measurable; patch generation combines deterministic and model-backed paths; the
learned decision gate removes the old low-risk magic threshold when enabled; safety and plausibility are
enforced through explicit gates; and the validation harness provides bounded CI-style re-execution. The
benchmark scripts and acceptance tests then exercise the same engine as both an experimental subject and a
usable tool.

With the implementation now grounded in code, the next chapter can describe how the system was evaluated:
which datasets were used, how baselines were defined, what metrics were frozen in advance, and how the project
maintained research discipline through phase-by-phase pre-registration.
